\documentclass[a4paper,12pt]{ctexart}
\usepackage{../../teaching-style}

\begin{document}
\pagestyle{empty}

\maintitle{线性代数初步}
\begin{center}
  \large 高学段 · 重结构轻计算，非工科数学
\end{center}
\vspace{0.3cm}

\chaptertitle{第一章\quad 线性方程组与矩阵}

\sectiontitle{1.1\quad 线性方程组的消元法}

\knowbox{
\textbf{线性方程组}：$\begin{cases}a_{11}x_1+a_{12}x_2+\cdots+a_{1n}x_n=b_1\\\vdots\\a_{m1}x_1+a_{m2}x_2+\cdots+a_{mn}x_n=b_m\end{cases}$\\
\textbf{高斯消元法}：通过行变换（交换两行、某行乘非零常数、某行倍数加到另一行）化为阶梯形。\\
\textbf{增广矩阵}：将系数和常数写成一个矩阵。\\
三种情况：唯一解、无穷多解、无解。}

\sectiontitle{1.2\quad 矩阵的概念与运算}

\knowbox{
矩阵$A_{m\times n}$：$m$行$n$列的数表。\\
加法：对应元素相加（同型矩阵）。\\
数乘：每个元素乘以该数。\\
乘法：$A_{m\times n}B_{n\times p}=C_{m\times p}$，$c_{ij}=\sum_{k=1}^n a_{ik}b_{kj}$（前行$\times$后列）。\\
注意：矩阵乘法不满足交换律（$AB\neq BA$一般）。}

\sectiontitle{1.3\quad 矩阵的秩与逆矩阵}

\knowbox{
\textbf{秩}：矩阵通过行变换化为阶梯形后非零行的行数，记作$r(A)$。\\
秩是矩阵最本质的不变量之一。\\
\textbf{逆矩阵}：若$AB=BA=I$，则$B=A^{-1}$。\\
可逆条件：$n$阶方阵$A$可逆$\Leftrightarrow r(A)=n\Leftrightarrow\det A\neq0$。\\
求逆：$(A|I)\xrightarrow{\text{行变换}}(I|A^{-1})$。}

\sectiontitle{习题1}

\begin{exercise}
\item 用高斯消元法解：$\begin{cases}2x+3y=7\\x-2y=-4\end{cases}$。
\item 计算矩阵乘积：$\begin{pmatrix}1&2\\3&4\end{pmatrix}\begin{pmatrix}0&1\\1&0\end{pmatrix}$。
\item 求$\begin{pmatrix}2&3&1\\0&1&-1\\0&0&2\end{pmatrix}$的秩。
\item 求$\begin{pmatrix}1&2\\3&4\end{pmatrix}$的逆矩阵。
\end{exercise}

\newpage

\chaptertitle{第二章\quad 行列式}

\sectiontitle{2.1\quad 行列式的定义与性质}

\knowbox{
二阶行列式：$\begin{vmatrix}a&b\\c&d\end{vmatrix}=ad-bc$。\\
三阶行列式（展开）：$\begin{vmatrix}a_{11}&a_{12}&a_{13}\\a_{21}&a_{22}&a_{23}\\a_{31}&a_{32}&a_{33}\end{vmatrix}=a_{11}a_{22}a_{33}+a_{12}a_{23}a_{31}+a_{13}a_{21}a_{32}-a_{13}a_{22}a_{31}-a_{12}a_{21}a_{33}-a_{11}a_{23}a_{32}$。\\
\textbf{性质}：\\
1. 转置行列式值不变。\\
2. 交换两行（列）变号。\\
3. 某行（列）乘$k$，值乘$k$。\\
4. 两行（列）成比例，值为$0$。\\
5. 某行（列）写成两数和，可拆成两个行列式之和。}

\sectiontitle{2.2\quad 克拉默法则}

\knowbox{
对于$n$元线性方程组$Ax=b$，若$\det A\neq0$，则唯一解为$x_i=\dfrac{\det A_i}{\det A}$，\\
其中$A_i$是将$A$的第$i$列换成$b$得到的矩阵。\\
理论意义重大，但实际计算中$n$较大时运算量巨大（高联二试可能涉及$n=2,3$的情况）。}

\sectiontitle{习题2}

\begin{exercise}
\item 计算：$\begin{vmatrix}2&3\\1&-4\end{vmatrix}$；$\begin{vmatrix}\sin\theta&-\cos\theta\\\cos\theta&\sin\theta\end{vmatrix}$。
\item 计算：$\begin{vmatrix}1&2&3\\2&3&1\\3&1&2\end{vmatrix}$。
\item 用克拉默法则解：$\begin{cases}2x+3y=7\\x-2y=-4\end{cases}$。
\end{exercise}

\newpage

\chaptertitle{第三章\quad 向量空间}

\sectiontitle{3.1\quad 向量与线性组合}

\knowbox{
向量$\vec v=(v_1,v_2,\dots,v_n)$是$n$元有序数组。\\
\textbf{线性组合}：$k_1\vec v_1+k_2\vec v_2+\cdots+k_m\vec v_m$。\\
\textbf{线性相关}：存在不全为零的$k_i$使线性组合$=0$。\\
\textbf{线性无关}：只有$k_i$全为零才能使线性组合$=0$。\\
几何意义（$\mathbb R^2$中）：两个向量线性相关$\Leftrightarrow$共线；三个向量一定线性相关。}

\sectiontitle{3.2\quad 基与维数}

\knowbox{
\textbf{基}：向量空间中一组线性无关且能生成整个空间的向量。\\
\textbf{维数}：基中向量的个数。\\
$\mathbb R^n$的标准基：$\vec e_1=(1,0,\dots,0)$，$\vec e_2=(0,1,\dots,0)$，$\dots$，$\vec e_n=(0,0,\dots,1)$。\\
$n$维向量空间中任意$n$个线性无关的向量构成一组基。\\
\textbf{坐标}：向量在给定基下的表示系数。\\
\textbf{基变换}：不同基之间的坐标变换通过过渡矩阵实现。}

\sectiontitle{习题3}

\begin{exercise}
\item 判断$\vec a=(1,2)$，$\vec b=(2,4)$是否线性相关。
\item 判断$\vec a=(1,0,1)$，$\vec b=(0,1,1)$，$\vec c=(1,-1,0)$是否线性无关。
\item 求$\vec v=(3,4)$在标准基$\{\vec e_1,\vec e_2\}$下的坐标。
\item 证明：$\mathbb R^2$中任意三个向量一定线性相关。
\end{exercise}

\newpage

\chaptertitle{第四章\quad 特征值与特征向量}

\sectiontitle{4.1\quad 特征值与特征向量的概念}

\knowbox{
对于$n$阶方阵$A$，若$A\vec v=\lambda\vec v$（$\vec v\neq\vec0$），则$\lambda$为\textbf{特征值}，$\vec v$为\textbf{特征向量}。\\
几何意义：$A$作用于$\vec v$相当于将$\vec v$伸缩$\lambda$倍（方向不变或相反）。\\
特征方程：$\det(A-\lambda I)=0$，展开得$\lambda$的$n$次多项式。}

\sectiontitle{4.2\quad 特征值的性质}

\knowbox{
1. 特征值之和$=$矩阵的迹（对角线元素之和）。\\
2. 特征值之积$=\det A$。\\
3. 不同特征值对应的特征向量线性无关。\\
4. 实对称矩阵的特征值都是实数，且可对角化。\\
\textbf{对角化}：若$A$有$n$个线性无关的特征向量，则$A=P^{-1}DP$，$D$为对角矩阵（元素为特征值）。\\
对角化的意义：简化计算$A^n$、解微分方程组等。}

\sectiontitle{习题4}

\begin{exercise}
\item 求$A=\begin{pmatrix}2&1\\1&2\end{pmatrix}$的特征值和特征向量。
\item 求$A=\begin{pmatrix}1&2\\3&2\end{pmatrix}$的特征值。
\item 若$A$的迹为$5$，$\det A=6$，求$A$的特征值。
\item 利用对角化计算$A^{10}$，其中$A=\begin{pmatrix}2&0\\0&3\end{pmatrix}$。
\end{exercise}

\newpage
\sectiontitle{参考答案}

\subsection*{习题1}
\begin{enumerate}
  \item $x=\dfrac27$，$y=\dfrac{15}{7}$
  \item $\begin{pmatrix}2&1\\4&3\end{pmatrix}$
  \item 秩为$3$
  \item $A^{-1}=\begin{pmatrix}-2&1\\\dfrac32&-\dfrac12\end{pmatrix}$
\end{enumerate}

\subsection*{习题2}
\begin{enumerate}
  \item $-11$；$1$
  \item $-18$
  \item $x=\dfrac27$，$y=\dfrac{15}{7}$
\end{enumerate}

\subsection*{习题3}
\begin{enumerate}
  \item 线性相关（$\vec b=2\vec a$）
  \item 线性相关
  \item $(3,4)$
  \item $\mathbb R^2$维数为$2$，任意$3$个向量必线性相关
\end{enumerate}

\subsection*{习题4}
\begin{enumerate}
  \item $\lambda_1=1$，特征向量$t(1,-1)$；$\lambda_2=3$，特征向量$t(1,1)$（$t\neq0$）
  \item $\lambda_1=4$，$\lambda_2=-1$
  \item $2$和$3$
  \item $A^{10}=\begin{pmatrix}2^{10}&0\\0&3^{10}\end{pmatrix}=\begin{pmatrix}1024&0\\0&59049\end{pmatrix}$
\end{enumerate}

\end{document}
